\section{System Model and Overview}
\label{sec:model}

\subsection{Participants, Outputs, and Certificates}
\label{sec:model-participants}
\label{sec:model-certificates}

Next operates a native UTXO ledger. CAL carries payment value, and FUEL pays processing fees. Users authorize payments that consume inputs and create recipient and change outputs. Registered organizations certify payments and hold additional CAL in public reserve accounts. Users retain ownership of their principal, although approval locks the selected inputs against conflicting spending. No bilateral channel is opened. Fees normally use final FUEL outputs owned by the payer; an organization may authorize sponsorship.

Each organization has four fixed members and a replaceable collector. Members independently check requests and record input locks and capacity reservations before signing. The collector distributes requests and assembles three matching signatures into a quorum certificate (QC), without acquiring spending authority. A TXCer combines this QC with the certified payment summary. A separate committee $\mathcal C$ of four equally weighted validators orders and executes public commands. It requires three votes and maintains outputs, reserves, and obligations.

An output specifies its asset, amount, owner, and consumption route. All principal inputs of a fast payment name one consumption organization, which certifies that payment and issues the certificate for its new outputs. Each new output's route identifies the organization that must certify a later payment spending it. The \emph{source} of an output (also its parent) is its creating transaction; a \emph{successor} (also its consumer or child) spends it. Child outputs are those created by the consumer. Each payment also carries an Intent, an identifier for one payer operation across organizations.

An output identifier $x$ binds the network, creating transaction, and output position. Spending is indexed by $(x,i)$, with $i=0$ for ordinary outputs. A TXCer binds the network, transaction identity, issuer, configuration, epoch, rules, output commitments and amounts, and admission grants. Its three distinct member signatures cover the same fact $c_T$. Output bodies may be delivered separately; recipients check commitment, position, and ownership before recording an output. Table~\ref{tab:model-notation} summarizes the notation.

\begin{table}[!b]
\caption{Principal notation.}
\label{tab:model-notation}
\centering
\begin{tabular}{@{}p{0.23\columnwidth}p{0.69\columnwidth}@{}}
\toprule
Symbol & Meaning \\
\midrule
$G,\mathcal C$ & Guarantor organization and public committee \\
$T,c_T$ & Payment and its certified summary fact \\
$(x,i)$ & Output identifier and spend instance \\
$g,B_g$ & Grant and cumulative authorization \\
$a_T,p_T,\mathsf{Rec}_T$ & Total CAL outputs, gross compensation, and recovered compensation \\
$\mathcal O_x,\mathcal D_x$ & Direct obligation for a missing output and its immutable compensation decision \\
$(H_b,P_b)$ & Stable block hash and original part-set header \\
$\rho_\ell,\rho_p$ & Committed logical and installed physical revisions \\
\bottomrule
\end{tabular}
\end{table}

\subsection{Fault and State Assumptions}
\label{sec:model-assumptions}

\noindent\textbf{Faults and cryptography.}
The adversary corrupts statically at most one member per organization and one committee validator throughout an execution; users and collectors behave arbitrarily. Honest participants share the network configuration, genesis, rules, and organization registry, and each route names one active configuration and one consumption-lock domain. Signatures are unforgeable, canonical digests and original block commitments are binding, and messages may be delayed arbitrarily.

\noindent\textbf{State.}
Honest signing state is atomic and never forgotten or rolled back. Before a member signs, its approval records the payment, the admission vector, and the direct input certificates of the request; the quorum certificate is not part of it. Public execution is deterministic and atomically committed, and followers apply authenticated consecutive blocks, updating state and cursor together. Compensation and representation read original blocks and their metadata, so correct committee replicas retain both at every height that holds an open obligation or a decided slot awaiting representation.

\noindent\textbf{Adaptation and progress.}
A trusted offline dealer generates a three-of-four adaptation key and distributes one share to each of the four public committee validators. The same static, at-most-one Byzantine assumption applies. Adaptation requires three valid contributions; three responsive honest validators suffice to provide them. Any adapted body must still be authorized by a committed command. Progress needs responsive honest participants, available evidence, sufficient fees and resources, and fair delivery and scheduling. The experimental in-memory configuration keeps state atomic only while the process runs, and each run starts from fresh genesis.

\subsection{Certified Continuation}
\label{sec:model-example}

Consider Alice's payment $T_A$, which spends a final 100-CAL input, creates a 40-CAL output $x$ for Bob, and returns 60~CAL as change. Its inputs name $G_A$, while $x$ names $G_B$. $G_A$ certifies $T_A$ and remains the issuer responsible for $x$. Bob checks the certificate against $x$'s body, commitment, position, ownership, and trusted configuration before recording his receipt.

Bob can then authorize $T_B$ to spend $x$ on a 40-CAL payment to Carol and request certification from $G_B$ before $T_A$ executes publicly. $G_B$ evaluates $T_B$ against its input locks, Intent records, fees, and available resources, and issues the certificate for $T_B$'s outputs if a quorum approves. $G_A$ and $G_B$ may be the same organization. The certificate gives Bob authenticated evidence for this request; a public obligation for a missing $x$ arises when $T_B$ successfully executes. Recipient verification and atomic local recording complete the fast path, while complete-material replication and public submission proceed in the background. Figure~\ref{fig:architecture} follows these two payments.

\begin{figure*}[t]
\centering
\begin{tikzpicture}[x=1cm,y=1cm,font=\footnotesize,
  box/.style={draw=black!55,rounded corners=2pt,align=center,minimum height=1.0cm,text width=2.6cm,fill=white},
  flow/.style={-{Stealth[length=2mm]},thick,draw=blue!60!black},
  bgflow/.style={-{Stealth[length=2mm]},draw=black!65,dashed},
  lab/.style={font=\scriptsize,align=center,fill=white,inner sep=2pt}]
\node[box,text width=1.9cm] (sender) at (0,0) {Sender wallet\\\textbf{Authorize payment}};
\node[box,text width=3.3cm,fill=blue!4] (org) at (4,0) {Consumption organization\\4 independent members\\\textbf{3 matching signatures}};
\node[box,fill=blue!4] (wallet) at (8.8,0) {Recipient wallet\\Verify and record TXCer\\\textbf{Ready to continue}};
\node[box] (next) at (13,0) {Next payment\\Same or another\\registered organization};
\draw[flow] (sender) -- node[lab,above]{Transaction} (org);
\draw[flow] (org) -- node[lab,above]{Output + TXCer} (wallet);
\draw[flow] (wallet) -- node[lab,above]{Direct evidence} (next);
\node[box,text width=3.3cm,fill=green!4] (comm) at (4,-3) {Public committee\\\textbf{Execute child payment}\\Consume once; record direct gap};
\node[box,fill=green!4] (resolve) at (8.8,-3) {Await source\\Arrival closes\\the obligation};
\node[box,fill=orange!6] (repair) at (13,-3) {Authorized repair\\Reserve debit + revision\\Late source: repay issuer\\Background installation};
\draw[bgflow] (org) -- node[lab,right]{Public submission} (comm);
\node[lab,left] at (2.1,-1.25) {Background INSTALL\\replicates complete evidence\\within the organization};
\draw[bgflow] (comm) -- node[lab,above]{Direct\\obligation} (resolve);
\draw[bgflow] (resolve) -- node[lab,above]{Still missing\\at deadline} (repair);
\node[align=center,font=\scriptsize,text=black!75] at (8.6,-4.5) {Source resolution preserves the authorized outputs and successor references.\\Wallets and members follow authenticated public blocks; no per-payment committee notification.};
\end{tikzpicture}
\caption{Continuous payment and public source resolution in Next. Solid arrows form the recipient's continuation path. Dashed arrows denote background work. Each subsequent payment is independently authorized; an unresolved source is assigned to its direct issuer rather than propagated as an ancestor-waiting requirement. Historical repair is authorized by a later public command; a source that executes after compensation repays the compensating issuer without crediting the recipient again.}
\label{fig:architecture}
\end{figure*}


\subsection{Public Consumption and Source Resolution}
\label{sec:model-contract}

If $T_A$ executes before $T_B$, $x$ is a final public output when $T_B$ consumes it, and no missing-source obligation is needed. If $T_B$ passes public checks and executes first, the committee atomically records the consumption of $x$, creates $T_B$'s final outputs, and registers a 40-CAL obligation against $G_A$. Carol can spend her output under the ordinary final-input rules while this obligation remains Open; her next payment still requires certification and admission by its designated organization. The block carrying $T_B$ also exposes $T_A$'s certificate. Original signers in $G_A$ can combine that certificate with their retained approval material and resubmit $T_A$ without new owner or member signatures.

If $T_A$ executes while the obligation is Open, it fulfills the obligation and closes the gap. This remains possible after the deadline if compensation has not executed. Once the deadline anchored in public block time is reached, an included valid compensation decision can instead debit $G_A$'s reserve by 40~CAL and close the gap. $T_B$'s principal effects remain unchanged. Closing each input obligation updates $T_B$'s fee accounting; its escrow closes after all its input obligations are resolved.

If $T_A$ executes after compensation, it consumes its own inputs and repays the same reserve account by 40~CAL instead of recreating $x$. Unconsumed change is created normally, and the compensation decision remains recorded. If $T_A$ never executes, the compensated 40~CAL remains $G_A$'s loss. If neither $T_A$ nor a valid successor spending $x$ executes, Bob's certificate provides no certificate-only redemption or unconditional exit; waiting alone does not create one. Table~\ref{tab:model-contract} states the guarantees by stage.

\begin{table}[!t]
\caption{Service contract by stage.}
\label{tab:model-contract}
\centering\small
\begin{tabular}{@{}>{\raggedright\arraybackslash}p{0.24\columnwidth}>{\raggedright\arraybackslash}p{0.69\columnwidth}@{}}
\toprule
Trigger & Established state \\
\midrule
Verified receipt &
The issuer's QC binds the output and its owner. The recipient may
request a successor, subject to its own admission checks. Receipt
alone creates no public compensation obligation. \\
Public consumption of a missing certified output &
Successful execution atomically consumes the input, registers its
direct issuer's obligation, and creates the child outputs.
Authorization, fees, consumption state, Intent, and coverage must pass the
public checks. Under the stated fault and state model, valid certificates have sufficient CAL authorization (Theorem~\ref{thm:security-coverage}). \\
Liability closure &
Source execution fulfills the obligation; otherwise an included,
valid post-deadline decision debits the issuer's CAL reserve.
A later source repays it without a second credit to the recipient.
The deadline alone does not guarantee completion. \\
\bottomrule
\end{tabular}
\end{table}

\subsection{Signing Capacity, Public Coverage, and Gaps}
\label{sec:model-layers}

Three quantities account for different stages of this example. Signing capacity limits each member's approvals. Public coverage tracks the outputs of a registered certificate. An open gap records the value of missing outputs already consumed publicly. The issuer's reserve account supplies the CAL used to close a gap through compensation.

Each member approving $T_A$ reserves 100~CAL, covering both $x$ and Alice's change. This reservation protects commitments that may circulate before public registration. When $T_B$ first consumes the missing $x$, the committee registers coverage for all 100~CAL of $T_A$'s outputs, but opens only a 40-CAL gap. The remaining 60~CAL stays covered. Registration leaves the reserve balance unchanged; compensation subsequently debits 40~CAL. Members restore their recorded CAL capacity after authenticating $T_A$'s successful execution; compensation alone leaves that capacity reserved.

A grant $g$ with cumulative authorization $B_g$ CAL gives each member $\lfloor2B_g/3\rfloor$ signing capacity, partitioned among its local execution Workers. Workers are internal processing units with resource slices, not additional validators. Each approving member reserves $a_T$, the full CAL output value including change. These overlapping signing capacities do not multiply the backing funds. Honest signers retain reservations for outstanding certificate liabilities, allowing the coverage proof to count certificates even when they remain unpublished. Genesis checks aggregate grants against their backing accounts; authorized top-ups increase both balance and grant, and authorizations never decrease.

Any client may collect approvals without spending or fee-sponsorship authority. A member atomically checks the Intent, input locks, fees, grants, and all applicable resource slices before approving. Insufficient capacity rejects that approval without committing changes or returning a signature. Retrying an existing, non-invalidated approval adds no reservation. New requests still need available capacity at a quorum. Partial or abandoned approvals may occupy capacity without producing a certificate, so new certification can stop while the reserve still has funds. Stoppage does not cancel certificates or unlock inputs.

Organizations keep local Intent records, while public execution accepts at most one transaction for the operation. Two transactions with the same Intent and disjoint inputs may therefore be certified by different organizations, although at most one can execute publicly. An unexecuted loser retains its signing reservation even after compensation. Exhaustion can stop other users' payments routed to that organization, including payments spending final outputs. FUEL expenditure and locked source inputs remain separate costs; fee sponsorship does not transfer the direct issuer's CAL liability to the fee payer.

The implementation has no payer-level CAL admission bound or cross-collector unfinished-request bound; fee-sponsorship policies and gateway concurrency limits do not provide either. Section~\ref{sec:security-composition} gives the workload and capacity conditions for continued admission and bounds the issuer's net unrepaid compensation.

\subsection{Economic Closure and Historical Records}
\label{sec:model-records}

After compensation, $T_B$'s original input still identifies $x$, the object accepted when $T_B$ executed. A local archival operator may need to report $G_A$'s reserve advance and export a record showing that advance at the original input position. A decision-authorized repair records the reserve debit in that input while preserving payment authorization, output references, and the complete original BlockID: the block-header hash together with the original part-set header.

At a verified and executed public prefix $H$, the reader uses one read transaction on its state $V_H$ and retained history to obtain the authorized representation, permanent decision, and current obligation status. After repayment, the reserve reference still records the historical advance. An ordinary decision index can answer the same economic questions; repair additionally places the debit reference inside a body compatible with the original complete BlockID. Logical revision and physical installation advance separately, and replay uses retained original commands. Repair requires a committed compensation decision; compensation and repayment proceed without waiting for adaptation or installation.

Table~\ref{tab:model-states} distinguishes these states. For example, an obligation can already be Compensated while its input still has \texttt{OriginalFunding} and no physical revision has been installed. Section~\ref{sec:protocol-reader} specifies the local reading and export interface.


\begin{table}[!t]
\caption{Independent state dimensions.}
\label{tab:model-states}
\centering\small
\setlength{\tabcolsep}{3pt}
\begin{tabular}{@{}p{0.24\columnwidth}p{0.70\columnwidth}@{}}
\toprule
Dimension & Meaning \\
\midrule
Public payment &
Successful execution (\texttt{Settled}) creates final outputs even with unresolved input obligations (\texttt{Pending}). When a delayed source executes, the value of each of its previously compensated outputs returns to the reserve account that paid the compensation; that output is not recreated. \\

Direct obligation &
Open$\to$Fulfilled, or Open$\to$Compensated$\to$Recovered.
Compensated denotes compensation; Recovered denotes subsequent repayment.
The compensation decision persists. \\

Logical representation &
The revision in $V_H$ determines the slot.
\texttt{OriginalFunding} does not imply absence of compensation.
\texttt{ReserveFunding} records a historical advance and remains
after repayment. \\

Physical installation &
$\rho_p$ records the installed block-store version.
Installation progress does not determine the logical answer at $H$
or change economic state. \\
\bottomrule
\end{tabular}
\end{table}

